\section{What to Learn: Information-Seeking Capabilities}
\label{sec:what_to_learn}

Effective information seeking requires more than retrieving relevant documents. A model must determine whether external information is needed, formulate information requests, explore available sources, evaluate retrieved evidence, integrate new information into its reasoning, and adapt subsequent actions accordingly. We organize these requirements into six interrelated capabilities. Rather than prescribing fixed pipeline stages, this taxonomy characterizes the behaviors that information-seeking policies are expected to acquire. We review representative approaches for each capability, while distinguishing learned behaviors from those enabled by inference-time mechanisms or external components.

\subsection{Epistemic Awareness and Retrieval Utility}
\label{sec:epistemic_awareness}

\textit{Epistemic awareness} concerns recognizing limitations of internal knowledge, while \textit{retrieval utility} concerns whether acquiring external information is expected to improve task performance. Early approaches primarily focus on retrieval necessity. \textit{When to Retrieve} learns when external retrieval should be invoked, while \textit{KBM} estimates knowledge boundaries to guide adaptive retrieval \cite{when_to_retrieve,kbm}. Related inference-time methods, including \textit{DRAGIN} and \textit{SeaKR}, use evolving information needs or uncertainty signals to trigger retrieval during generation \cite{dragin,seakr}. However, \textit{Adaptive Retrieval Without Self-Knowledge?} provides an important counterpoint: uncertainty or self-knowledge proxies do not necessarily identify whether retrieval will actually improve the answer \cite{adaptive_retrieval_without_self_knowledge}. These approaches therefore demonstrate selective retrieval behavior, but not yet calibrated prediction of retrieval value.

More recent research shifts attention toward the utility of information acquisition. Studies of LLM-specific retrieval utility and cost-aware evidence selection suggest that the usefulness of retrieved content depends on the consuming model, evidence quality, and acquisition cost \cite{llm_specific_utility,when_knowledge_not_free}. Building on this perspective, \textit{s3} trains a search policy with a downstream gain-over-RAG signal, while \textit{CoBRA} learns tool-use boundaries through counterfactual comparisons between answering with and without external tools \cite{s3,cobra}. \textit{Calibrate-Then-Act} further studies cost-aware action selection under an explicit calibrated answerability prior \cite{calibrate_then_act}. The latter is an important cost-sensitive decision baseline, but its calibrated prior is an external component and should not be interpreted directly as intrinsic calibration learned by the LLM. Other approaches optimize search behavior through information-gain-oriented rewards or reinforce the coordination of internal and external knowledge \cite{info_reasoner,ikea}.

Nevertheless, retrieval necessity and information value remain distinct learning targets. In particular, improvements obtained from realized retrieval outcomes do not necessarily establish that a policy can accurately predict retrieval benefits \emph{before} information is acquired. This distinction motivates separating learned retrieval routing, ex-post utility optimization, and calibrated, cost-sensitive information-acquisition decisions.

\subsection{Information-Need Formulation}
\label{sec:information_need_formulation}

\textit{Information-need formulation} concerns translating unresolved knowledge gaps into effective queries, subquestions, or retrieval plans. Early work investigates query rewriting as a learnable component of retrieval-augmented generation \cite{query_rewriting}. More recent approaches directly optimize query-generation policies using retrieval feedback. \textit{DeepRetrieval} learns queries adapted to search engines and retrievers, while query-expansion methods and \textit{ParallelSearch} extend formulation beyond a single request by generating complementary queries or decomposing complex information needs \cite{deepretrieval,expandsearch,parallelsearch}.

For iterative reasoning, information needs may evolve as new evidence becomes available. \textit{StepSearch} optimizes intermediate search behavior, whereas \textit{QREF} aligns query refinement with reasoning requirements and evidence coverage \cite{stepsearch,qref}. Together, these methods illustrate a progression from static reformulation toward reasoning-conditioned query policies.

However, retrieval effectiveness does not necessarily imply faithful or transferable formulation. \textit{Environment-aware IR} shows that learned query patterns can depend strongly on the retrieval backend, with potential limitations when transferred across retrievers \cite{environment_aware_ir}. We therefore distinguish three aspects of this capability: preserving the semantic information need, adapting queries to a particular retrieval environment, and transferring formulation behavior to unseen environments.

\subsection{Information-Space Exploration}
\label{sec:information_space_exploration}

\textit{Information-space exploration} concerns discovering useful information sources and previously uncovered evidence, rather than simply generating additional queries. \textit{MARAG-R1} investigates learned coordination across multiple retrieval tools, while diversity-driven policy optimization encourages broader exploration of search trajectories \cite{marag_r1,ddpo_search}. Complementary studies, including \textit{InfoMosaic-Bench} and \textit{Table-as-Search}, provide settings for evaluating multi-source acquisition and representing unresolved information gaps \cite{infomosaic_bench,table_as_search}. \textit{Agentic Search in the Wild} complements these controlled settings with large-scale analysis of real search trajectories, including exploration breadth, repetition, and evidence-conditioned query evolution \cite{agentic_search_in_the_wild}.

Importantly, query diversity does not necessarily translate into evidence diversity, and neither guarantees useful information coverage. Likewise, selecting a valid tool differs from identifying a source containing necessary evidence. Exploration is therefore best understood as a cross-cutting capability involving source selection, coverage assessment, and the discovery of information not yet available to the agent. The real-world trajectory evidence is useful for measuring these behaviors, but lexical novelty or trajectory diversity should not be treated as direct evidence of useful information gain.

\subsection{Evidence Evaluation}
\label{sec:evidence_evaluation}

\textit{Evidence evaluation} concerns determining whether acquired information is relevant, supportive, reliable, and sufficient for the current task. These judgments are related but distinct: a relevant passage may not support a particular claim, while a supporting passage may originate from an unreliable source.

\textbf{Relevance and support assessment} are addressed by methods that explicitly evaluate retrieved content. \textit{CRAG} introduces retrieval-quality assessment to guide corrective processing, while \textit{RLSeek} learns evidence-grounded judgments for hallucination detection \cite{crag,rlseek}. \textit{EAPO} further introduces evidence-quality rewards into policy optimization for long-context reasoning, but its evidence reward remains coupled to answer and reward-model signals \cite{eapo}. Evaluation frameworks such as \textit{RAGTruth} and \textit{RAGChecker} further characterize unsupported generation and retrieval--generation failures \cite{ragtruth,ragchecker}. Beyond individual passages, \textbf{evidence sufficiency} requires determining whether the available information supports the complete inference. \textit{S2G-RAG} explicitly models sufficiency and unresolved information gaps, while \textit{Sufficient Context} distinguishes insufficient retrieval context from failures to use available evidence \cite{s2g_rag,sufficient_context}.

A separate challenge concerns \textbf{source reliability and conflicting evidence}. \textit{AuthorityBench} investigates models' perception of source authority, whereas studies of conflicting evidence in fact-checking and web references examine how models respond to incompatible claims \cite{authoritybench,confact,conflicting_web_references}. Related counterfactual analyses suggest that evidence preferences can also be influenced by perceived relevance and presentation-related cues \cite{what_evidence_convincing}. \textit{(D)RAGged Into a Conflict} further shows that conflicts have different types and expected response behaviors, rather than forming a single binary condition \cite{dragged_into_conflict}. These findings highlight the difference between identifying apparently convincing evidence, classifying a conflict, and determining whether evidence is trustworthy.

Existing approaches make relevance, support, and sufficiency judgments increasingly explicit, with some directly optimized through training. However, detecting supporting content does not establish factual reliability, and recognizing disagreement does not necessarily imply resolving it correctly. These distinctions are essential when interpreting evidence-evaluation performance as a learned capability.

\subsection{Evidence Integration and Belief Revision}
\label{sec:evidence_integration}

\textit{Evidence integration} concerns whether acquired information appropriately influences a model's reasoning, conclusions, and subsequent information state. Unlike evidence evaluation, which assesses the quality of observations, integration requires using those observations to revise existing beliefs and construct evidence-grounded inferences. We distinguish three closely related aspects.

\textbf{Knowledge integration and belief revision.}
When retrieved information contradicts parametric knowledge, models must determine whether the external evidence warrants revising an existing conclusion. \textit{ClashEval} investigates conflicts between internal priors and external evidence, while \textit{Belief-R} probes whether models update or maintain beliefs appropriately when new premises are introduced \cite{clasheval,belief_r}. Studies of source balancing examine how models respond to assertions from users and documents, while \textit{Efficient Prior-Guided Reasoning} studies robust reasoning under noisy or conflicting retrieval contexts \cite{source_balancing,efficient_prior_guided_reasoning}. Learning-oriented approaches such as \textit{GuarantRAG} investigate stronger integration of retrieved knowledge, while \textit{FaithfulRAG} addresses fact-level conflicts between internal knowledge and external context \cite{guarantrag,faithfulrag}. These approaches highlight a central requirement: reliable belief revision should incorporate well-supported evidence without indiscriminately accepting potentially misleading external assertions.

\textbf{Multi-evidence composition.}
Complex information needs frequently require combining complementary facts distributed across documents. \textit{RAFT} trains models to utilize relevant evidence in the presence of distractors, while \textit{BRIEF} learns compressed evidence representations for multi-hop reasoning \cite{raft,brief}. \textit{TRACE the Evidence} explicitly constructs document-grounded multi-hop chains and provides a structured comparison between chain construction and all-document contexts \cite{trace_evidence}. These approaches facilitate evidence-conditioned inference, although improvements from context organization, compression, or chain construction do not necessarily establish that the reader has learned to combine every necessary premise.

\textbf{Faithful evidence utilization.}
Beyond generating answers consistent with retrieved context, models should produce conclusions that genuinely depend on the acquired information. \textit{CARE} explicitly encourages evidence-grounded reasoning, while \textit{MIRAGE} provides attribution-based analysis of retrieved evidence contributions \cite{care,mirage}. \textit{FRANQ} further distinguishes factuality from faithfulness: an answer may be factually correct without being supported by the current retrieval context, or faithful to an erroneous retrieved claim \cite{franq}. These studies motivate distinguishing \emph{output reliance}, where answers change with evidence; \emph{compositional reliance}, where multiple necessary facts are utilized; and \emph{process-level reliance}, where those facts causally influence the decision. Evidence of output sensitivity, factual correctness, or plausible attribution alone does not establish the strongest form of reliance. Attribution shortcuts caused by source metadata are treated separately in Section~5.

\textbf{Cross-turn information persistence.}
Evidence integration also extends to retaining and revising information across interactions. \textit{LongMemEval} evaluates long-term information use, while \textit{MemoryAgentBench} introduces incremental interaction, contradiction, temporal-update, and selective-forgetting tests \cite{longmemeval,memoryagentbench}. \textit{Memory-R1} learns memory-management operations and \textit{MemRL} adapts external episodic-memory utilities \cite{memory_r1,memrl}. These approaches provide complementary perspectives on persistent information use, although improvements from memory infrastructure should be distinguished from capabilities learned by the underlying language-model policy.

\subsection{Adaptive Search Control}
\label{sec:adaptive_search_control}

\textit{Adaptive search control} governs how an agent responds to information obtained during an ongoing search trajectory. Depending on the observed evidence, the model may continue searching, reformulate its query, verify an intermediate conclusion, switch information sources, or terminate the process.

Initial work on learned retrieval--reasoning interactions, including \textit{Self-RAG}, establishes retrieval and critique as components of model behavior \cite{selfrag}. More recent approaches such as \textit{Search-R1} and \textit{R1-Searcher} optimize iterative search and reasoning trajectories, while \textit{ZeroSearch} investigates learning search behavior through simulated interactions \cite{search_r1,r1_searcher,zerosearch}. These studies demonstrate that multi-step search behavior can be incorporated into learned policies.

Other approaches target more specific control decisions. \textit{AutoSearch} learns adaptive search depth, while \textit{One Policy, Any Budget} investigates budget-aware decisions across reasoning, searching, and answering \cite{autosearch,anysearch}. Research on beneficial search behaviors further identifies verification, adaptive search, and error recovery as important components of successful trajectories \cite{beneficial_search_behaviors}. \textit{CriticSearch} studies retrospective turn-level credit assignment, while \textit{CW-GRPO} reweights optimization according to the estimated contribution of retrieval actions to final reasoning success \cite{criticsearch,cw_grpo}.

Although these studies demonstrate improvements in iterative search and resource allocation, executing effective multi-step trajectories is not necessarily equivalent to responding appropriately to new observations. In particular, budget-sensitive stopping, hindsight credit assignment, and feedback-conditioned adaptation represent distinct behaviors. The defining capability is whether the agent can revise its subsequent actions according to the information actually obtained, rather than following a largely predetermined search pattern.
